% Run GCT_KRON_20261120; GPT-6 Astra Pro; effort not exposed.
\section{Uniform injectivity of the symmetric quadratic block}

\subsection*{QA033-SYM: one rectangular gate is settled in every size}
The symmetric map \(F\) in (CB.9) is injective for a nonempty
Zariski-open set of the distinct parameters \(a_1,\ldots,a_m\),
for every \(m\geq2\). The proof gives an explicit ordered-parameter
limit in which this rectangular map is injective. It does not
prove the mixed surjectivity or the final charge-map assertion
in (CB.14).

\paragraph{The ordered limit and its exact tensor.}
Set \(a_i=\epsilon^{-i}\), where \(\epsilon\) is a nonzero
parameter, and write \(\epsilon\to0\) for the coefficientwise
limit of the rational functions that follow. Let
\(\epsilon_i\) denote the standard basis of \(\mathbb C^m\).
The residue matrix of the binary operator \(1+t\partial_t\)
is
\[
 L_{ii}=1-a_i\sum_{j\ne i}\frac1{a_j-a_i},\qquad
 L_{ij}=-\frac{a_i}{a_j-a_i}\quad(i\ne j).
                                                               \tag{SI.1}
\]
Indeed, writing a binary polynomial of degree at most \(m-1\)
as \(\sum_j u_j\Delta_j\), evaluation of its derivative at
\(t=-a_i\), divided by \(\Delta_i(-a_i)\), gives (SI.1).
The same formula and (O-031.2) give \(e=-2Lp\).

The matrix \(L\) tends to the lower triangular matrix
\[
 (L_\infty)_{ii}=i,\qquad
 (L_\infty)_{ij}=1\ (j<i),\qquad
 (L_\infty)_{ij}=0\ (j>i).
                                                               \tag{SI.2}
\]
On the residue hyperplane, use the simple-root basis
\(f_i=\epsilon_i-\epsilon_{i+1}\), \(1\leq i<m\).
Then \(L_\infty f_i=i f_i\).

For an edge \(a<b\), multiply its symmetric gradient column
by \(a_b-a_a\). This is an invertible source scaling for every
nonzero \(\epsilon\) with distinct parameters. Since
\[
 \operatorname{Res}\frac{\Delta_{ab}}{\Delta}
                    =\frac{\epsilon_a-\epsilon_b}{a_b-a_a},
\]
the first line of (O-031.3) has the following finite limit:
\[
 -(p_a-p_b)L_\infty(\epsilon_a-\epsilon_b)
       +2\left(a p_a+\sum_{j<a}p_j\right)(\epsilon_a-\epsilon_b).
                                                               \tag{SI.3}
\]
Here the integer \(a\) in the parentheses is the vertex index.
In fact, \((a_a+a_b)/(a_b-a_a)\to1\),
\(a_b/(a_b-a_a)\to1\), and
\(a_a/(a_b-a_a)\to0\); these limits prove (SI.3) directly.

Consequently, in the output basis \(f_i\), the limiting
symmetric edge gradient is supported on \(a\leq i<b\), with
coefficient
\[
 g_{ab,i}=(2a-i)p_a+i p_b+2\sum_{j<a}p_j.
                                                               \tag{SI.4}
\]
Let \(F_m^\infty\) multiply these gradients by arbitrary
linear forms in \(p_1,\ldots,p_m\). We prove below that
\[
 F_m^\infty:\mathbb C^{\binom m2}\otimes
              \langle p_1,\ldots,p_m\rangle
       \longrightarrow
          \mathbb C^{m-1}\otimes
                 \operatorname{Sym}^2\langle p_1,\ldots,p_m\rangle
                   \quad\hbox{is injective.}
                                                               \tag{SI.5}
\]

\paragraph{Induction and the last output equation.}
A kernel vector is a family of linear forms \(X_{ab}(p)\),
one for every \(a<b\), such that
\[
 \sum_{a\leq i<b}g_{ab,i}X_{ab}=0
                      \qquad(1\leq i<m).
                                                               \tag{SI.6}
\]
For \(m=2\), this is \((p_1+p_2)X_{12}=0\), so the assertion
holds. Assume it for \(m-1\), and put \(k=m-1\).
The last output \(i=k\) involves only edges \(\{a,m\}\),
and gives
\[
 \sum_{a=1}^k\ell_a X_{am}=0,
 \qquad
 \ell_a=k p_m+(2a-k)p_a+2\sum_{j<a}p_j.
                                                               \tag{SI.7}
\]
These \(k\) linear forms are independent.
To see this, let \(L_k\) be the lower triangular matrix
with diagonal entries \(2a-k\) and strict lower entries \(2\).
The row coefficient matrix of the \(\ell_a\) is
\([L_k\mid k\mathbf1]\).
If \(k\) is odd, \(L_k\) is invertible.
If \(k=2h\) is even, its left kernel is one dimensional,
spanned by the vector with its first \(h\) entries equal
to \(1\) and its remaining entries zero: this follows
by solving the columns from \(k\) down to \(1\).
That vector has nonzero pairing \(h\) with \(\mathbf1\).
The augmented matrix therefore has rank \(k\) in both cases.

Every linear syzygy among independent linear forms is skew:
complete the \(\ell_a\) to a basis of the linear forms and
compare all quadratic coefficients in (SI.7). Terms involving
the extra basis vector vanish, and the remaining coefficient
matrix is skew-symmetric. Thus there is a constant skew matrix
\(T=(T_{ab})_{1\leq a,b\leq k}\) such that
\[
                     X_{am}=\sum_b T_{ab}\ell_b.
                                                               \tag{SI.8}
\]

For \(i<k\), the coefficient of \(p_m^2\) in (SI.6) is
\(ik\sum_{a\leq i}(T\mathbf1)_a\).
Successive values of \(i\), and then
\(\mathbf1^T T\mathbf1=0\), show that
\[
                     T\mathbf1=0,
                     \qquad\mathbf1^T T=0.
                                                               \tag{SI.9}
\]
In particular \(X_{am}\) has no \(p_m\) term. With
\(p^{\rm old}=(p_1,\ldots,p_k)^T\), set
\[
                  W=T L_k,\qquad
                  (X_{am})_{a=1}^k=W p^{\rm old}.
                                                               \tag{SI.10}
\]
All column sums of \(W\) are zero by (SI.9).

\paragraph{The coefficients containing \(p_m\).}
For an old edge \(a<b\leq k\), write
\(X_{ab}=Y_{ab}+c_{ab}p_m\), where \(Y_{ab}\) involves
only the old variables. The coefficient of \(p_m\) in
(SI.6), as a linear form in those variables, is
\[
 \sum_{a\leq i<b\leq k}g_{ab,i}c_{ab}
                  +i\sum_{a\leq i}(Wp^{\rm old})_a=0
                          \qquad(1\leq i<k).
                                                               \tag{SI.11}
\]
For \(j>i\), comparison of the \(p_j\) coefficient yields
\(\sum_{a\leq i}(c_{aj}+W_{aj})=0\).
Subtracting consecutive such equations gives
\[
                         c_{ab}=-W_{ab}\qquad(a<b\leq k).
                                                               \tag{SI.12}
\]
For \(j\leq i\), substitute (SI.12) into the corresponding
coefficient of (SI.11). The result is
\[
 i\sum_{a\leq i}W_{aj}
 -(2j-i)\sum_{b>i}W_{jb}
 -2\sum_{a=j+1}^i\sum_{b>i}W_{ab}=0.
                                                               \tag{SI.13}
\]
These equations force \(T=0\), as follows.

\paragraph{Descending-column elimination with nonzero coefficients.}
Suppose that all rows and columns of \(T\) with index
greater than \(l\) are zero; this is vacuous for \(l=k\).
Take \(i=l-1\) in (SI.13).
The zero column sums of \(W\), and the vanishing rows of
\(W\) beyond \(l\), give
\[
 -(l-1)W_{lj}-(2j-l+1)W_{jl}
             -2\sum_{a=j+1}^{l-1}W_{al}=0
                              \qquad(1\leq j<l).
                                                               \tag{SI.14}
\]
Let \(h=k-l\geq0\). From \(W=T L_k\), skew symmetry,
and the same vanishing rows and columns,
\[
 W_{jl}=(2l-k)T_{jl},\qquad
 W_{lj}=-(2j-k)T_{jl}
                       -2\sum_{a=j+1}^{l-1}T_{al}.
                                                               \tag{SI.15}
\]
Substitution into (SI.14), followed by division by \(2\),
reduces it exactly to
\[
 -\bigl(j+h(l-1-j)\bigr)T_{jl}
                    +(h-1)\sum_{a=j+1}^{l-1}T_{al}=0.
                                                               \tag{SI.16}
\]
Start with \(j=l-1\), and then decrease \(j\).
The sum in (SI.16) is zero at each step because its entries
have already vanished. The coefficient
\(j+h(l-1-j)\) is a positive integer for every
\(1\leq j<l\). Thus the entire \(l\)-th column, and by
skew symmetry the \(l\)-th row, vanish.
Decreasing \(l\) from \(k\) to \(2\) proves \(T=0\).

Equations (SI.10) and (SI.12) now imply
\(X_{am}=0\) and \(c_{ab}=0\). The remaining equations
(SI.6), for \(i<k\), are exactly the kernel equations
for \(F_{m-1}^\infty\) on the \(Y_{ab}\), because
the coefficients (SI.4) do not depend on \(m\).
The induction hypothesis kills all \(Y_{ab}\).
This proves (SI.5) for every \(m\geq2\).

\paragraph{Transfer to the actual rational residue tensor.}
An injective rectangular matrix has a nonzero maximal
minor. The entries of the scaled residue matrix used
above are rational functions of \(\epsilon\) regular
at zero, and such a minor has nonzero limiting value
by (SI.5). It is therefore a nonzero rational function.
For generic nonzero \(\epsilon\), the original symmetric
map has full column rank, because the edge scalings
are invertible. Equivalently, its maximal-rank locus
in the space of distinct parameters \((a_i)\) is a
nonempty Zariski-open set.

This settles the injectivity condition on \(F\) in
(CB.14) uniformly. The mixed map \(A\) and the induced
charge map remain independent gates; injectivity alone
does not imply the arbitrary-size endpoint theorem.
The ordered-limit computations in
\path{agents/jet_review/check_residue_ordered_limit_blocks.py}
were used to identify this induction, but the proof
above does not rely on finite rank tests.
The coefficient-family and comodule dictionary is exactly
the one in QA033-CB and O-031: SGA 3, I, Propositions
4.7.2--4.7.3, pp.~34--36. The proof concerns an actual
permanent coefficient differential; its generic-rank step
is the elementary nonvanishing of a polynomial minor.
